% !TeX root = ../papers/nc-p3.tex
\section{混合型の正規形、点の有限性と半周期成分への特殊化}
\label{sec:mixed-point-orbits}
\begin{corollary}\label{cor:all-defective-one-sym-finite}
片側対称な標準有限Frobenius窓の四辺が有限で、反対階数が9以下なら、
長さ1の点がなくても全次数の正則性は有限に決まる。
対称核に階数2以下の非零元があれば正則である。それ以外は両対称な正則窓に限る。
\end{corollary}
\begin{proof}
低階数元から中支持3を得る定理\ref{thm:rank-nine-finite-symmetry}の末尾の証明は核の次元を使わない。
それ以外では任意の非零核元が階数3または4である。階数3では定理\ref{thm:cone-kernel-binary-transfer}、
階数4では定理\ref{thm:smooth-symmetric-kernel-all}の証明を用いる。
後者で使う核の条件は、選んだ行列式の消滅、すなわち二組の二項位置の各対称性だけである。
いずれも両対称または指数成長となり、前者の判定は定理\ref{thm:bisymmetric-two-length-test}である。
\end{proof}

\begin{proposition}\label{prop:mixed-quartic-normal-forms}
前節の混合型で正則な全パラメータは、点グラフの次元によらず反対階数10、最小周期2である。
印付き正規形の可換環は $4=n_1+\cdots+n_l$ に対する
$\prod_i k[t_i]/t_i^{n_i}$ または $k[u,v]/(u^2,v^2)$ である。
\begin{enumerate}
\item 前者では各局所単位の値 $\omega_i$ の和は1であり、$n_i=1$ のとき $\omega_i\ne0$ である。
$n_i>1$ なら $\tau(t_i^{n_i-1})=1$、その間の冪の値は全て零とする。
残余同型は同じ $(n_i,\omega_i)$ の置換と $t_i\mapsto\zeta_i t_i$, $\zeta_i^{n_i-1}=1$ ($n_i>1$) で尽きる。
\item 後者では $\tau(u)=\tau(v)=0$, $\tau(uv)=1$ とする。
残余群は $(u,v)\mapsto(cu,c^{-1}v)$ と $u,v$ の交換が生成する。
\end{enumerate}
いずれも $\ker\tau$ 上の射影的な可逆対称三行列をこの残余群の合同作用で比較すれば必要十分である。
\end{proposition}
\begin{proof}
非被約な局所Frobenius環で接空間が一次元なら一文字の打切り環である。
接空間が二次元以上ならsocleの次元1と長さ4からHilbert関数は $1,2,1$ で、
中央二次元の積はsocleに値を持つ非退化対称形式となる。
二つの等方方向を取れば第二の環を得る。三次元接空間は平方零の三次元socleとなり不可能である。
一文字の環では最上冪の値は非零である。文字の倍率でそれを1にしてから
$t\mapsto t+c_jt^j$ ($j=2,\ldots,n-1$) で順に $\tau(t^{n-j})=0$ とする。
実際、変化量は $(n-j)c_j$ で既に揃えた高次の値は保たれる。
従ってこの表示の自己同型は表示した根の倍率だけである。
第二の場合は $uv$ の倍率を揃え、$u,v$ へsocle元を加えて各値を零とする。
自己同型は二つの等方方向を保ち、跡の条件がsocleへの追加を禁じる。

階数について、前節の残余対合が階数1となるなら
$Y=(L^t L)^{-1}$ は $I,Q_b$ 全てと跡直交する。
$U$ の直交基底 $e_i$ を使い $a=\sum Y_{ij} e_i e_j\in K$ とする。
$\tau(a)=-\operatorname{tr}Y=0$ であり、$a$ の各 $z\in U$ との積跡は
$-\operatorname{tr}(Y B_z)=0$ である。従って $a=0$ となる。
一つの局所因子の剰余体で評価し、その指標を $\chi$、socleの生成元を $d$ とする。
$e_i\cdot(d-\tau(d)1)=-\tau(d)\chi(e_i)$ なので
$y_d\cdot y_d=\tau(d)^2\chi(a)=0$ である。正則性条件に反する。
残余対合の階数0も既に排除したので全正則点の反対階数は10である。
周期1なら二つの反対位置で対称となり、両対称な標準AS窓の階数制限に反する。
\end{proof}

\begin{proposition}\label{prop:mixed-finite-eigenvectors}
正則な混合型で四辺が有限となる必要十分条件は、$K$ が一文字局所因子の積であり、
写像 $x\mapsto B_{(L^tL)x}x$ の $\mathbb P U$ 上の固有方向が有限となることである。
後者は $x\wedge B_{(L^tL)x}x$ の三つの三次式に非定数共通因子がないことと同値である
（三式全零の場合を除く）。
\end{proposition}
\begin{proof}
二周期の二つの辺は補文字 $r$ の符号変更と順序反転で対応するので、第一の辺の台を調べればよい。
原点対の外では $x,y$ はともに非零である。$r=0$ なら $y$ は $Lx$ の倍なので
$s'x+B_{n_y}x=0$ はちょうど固有方向の条件になる。
$r\ne0$ なら $y\cdot y=(Lx)\cdot y=0$ が外積式から従う。後者により
$x(s'+n_y)=0$ が $K$ の全積として成り立つ。
逆に $k'=s'+n_y$ が $x\in U$ を消し $y\cdot y=0$ なら $(Lx)\cdot y=0$ であり、
三次元の等方直線の性質から $(Lx)\times y$ は $y$ の倍で、一意な $r$ が取れる。
この逆構成で $r=0$ も許す。

$\operatorname{ann}(k')$ は $U$ と非自明に交わるので次元2以上である。
一次元イデアルは $U$ と交わらないからである。
一方、$k'$ が一次元イデアルを生成するなら $y$ は前節のsocle判定で非等方となる。
従ってこの構成では $I=k'K$ はちょうど二次元、$kx=\operatorname{ann}(I)\cap U$ である。
各二次元イデアル $I$ の $F$ への持上げの行列鉛筆は純四元数二平面であり、
等方根は重複込み二個である。各根が実際に $I$ を生成することも前節で示した。
異なる $I$ は異なる奇点を与える。一文字局所環の積に二次元イデアルは有限個しかなく、
$k[u,v]/(u^2,v^2)$ には $\langle a u+b v,uv\rangle$ の射影直線がある。
これらと $r=0$ の全点で被覆されるので結論を得る。
\end{proof}

\begin{theorem}\label{thm:mixed-half-period-specialization}
全ての正則な混合型の窓は、九次元半周期成分の正則な枠の特殊化である。
有限点を持つ混合型の稠密部分は印を忘れて八次元の既約な軌跡をなし、
その閉包はこの成分で余次元1である。
\end{theorem}
\begin{proof}
各 $(K,\tau)$ を被約な長さ4の可換代数へ平滑化する。一文字局所因子は
$t^n$ を $\prod_{j=1}^n(t-ja)$ に、第二の環は $(u^2,v^2)$ を $(u(u-a),v(v-a))$
に変えればよい。パラメータは $a$ で、単項式基底は一定である。
$\tau$ の基底上の値を固定すれば完全対合は零ファイバーで非退化なのでその近傍でも非退化である。
特に一般被約ファイバーでは全ての原始冪等元で値が非零である。
$\ker\tau$ とその内積は束をなし、$L$ は基底上で延ばせる。随伴は延ばした内積で取ればよい。

この近傍上で $D=M_2\oplus K$ の四次元部分空間 $E'$ をGrassmann束全体に動かし、
$F'=(E')^\perp$ とする。二次積 $E'F',F'E'$ がともに $H$ となる条件は開であり、
指定した混合型点で満たす。両積を $H/kz$ に写した六成分を二周期の関係として取れば、
この点の元の関係へ特殊化する。四階テンソル自体も
$\ell(e f e' f')$ で規格化できる。これは $z=(1,1)$ を割った
相補二次積の対合であり、$H^\perp=kz$ から隣接階数6である。

一般点では $K=k^4$ に分裂し、両四平面の両環因子への射影は同型である。
そのとき偶文字を行列 $X$、奇文字を $\operatorname{adj}Y$ により四次元ずつ同定する。
$\omega_i=\tau(e_i)$ を用いスカラー積の四成分をそれぞれ $\omega_i$ 倍して読むと、
偶側の四座標を $u_i$、奇側を $v_i$（この倍率は片方に入れる）として
$\sum u_i\otimes v_i$ は行列式の偏極となり、$z$ は $(\omega_1,\ldots,\omega_4,I_2)$ となる。
従って定理\ref{thm:half-period-frame}の枠表示である。一般枠では各単一座標の行列式と
その双対側の行列式が非零なので、定理\ref{thm:half-period-semisimple}により正則である。
一般の枠が現れることは、$E'$ が二つの四次元の間の任意の同型のグラフを走り、
$F'$ は積跡直交で一意に定まるためである。必要なら $\tau(1)=1$ を保って $\tau$ も動かせる。

混合型の被約 $K$ では三つの重み自由度と $U$ 上の射影二次形式の五つの自由度があり、
残余群は有限である。この八次元既約パラメータの一般点は正則で四辺が有限である。
実際、$K=k^4$ に等しい重みを取り、$(\mathbb Z/2)^2$ の三つの非自明指標に
$\sqrt{-1}$ を掛けた直交基底を $U$ に取って $L=I$ とすれば、$B_yx=\sqrt{-1}D(x,y)$
（ここで $D$ は定理\ref{thm:twenty-point-distinct-companion}の対称積）である。
補座標を定数倍すれば第一辺はその被約二十点の例になる。有限性・横断性と前節の非等方性は開条件である。
印の選び方が有限なので印なしの像も八次元であり、九次元成分内の閉包の余次元を得る。
非被約入力も同じ八次元混合型の閉包に入る。実際、先の平滑化で混合型表示を保って
形式を一般に動かすと一般ファイバーは正則かつ有限である。
\end{proof}
